\section*{Hoofdstuk \ref{ch:g1:watching-the-weather} --- Kijken naar het weer}

\begin{solution}{exo:g1:watching-the-weather:1}
Zonnig: \omwx{sun}\quad regen: \omwx{rain}\quad sneeuw: \omwx{snow}\quad
bewolkt: \omwx{cloud}. Elke duidelijke tekening van een zon, een wolk met
druppels, een wolk met sneeuwvlokken en een wolk alleen is goed.
\end{solution}

\begin{solution}{exo:g1:watching-the-weather:2}
Het park: zonnig en warm. De straat: sneeuw en kou. De heuvel: winderig
(en een paar wolken). Het hek van de school: mist.
\end{solution}

\begin{solution}{exo:g1:watching-the-weather:3}
Dinsdagochtend: bewolkt. Vrijdagmiddag: buien. Zondagmiddag: regen.
\end{solution}

\begin{solution}{exo:g1:watching-the-weather:4}
Woensdag: zowel 's ochtends als 's middags is er een zon getekend. (Op
donderdagmiddag staan er zon en wolken.)
\end{solution}

\begin{solution}{exo:g1:watching-the-weather:5}
De zonnehoed past bij warm en zonnig weer; laarzen en een regenjas bij
regen; een sjaal en handschoenen bij koud weer met sneeuw.
\end{solution}

\begin{solution}{exo:g1:watching-the-weather:6}
Winderig weer: de wind duwt tegen de vlag.
\end{solution}

\begin{solution}{exo:g1:watching-the-weather:7}
\begin{enumerate}
  \item Niet waar: het \omterm{def:g1:watching-the-weather:weather}{weer} kan in de loop van de dag veranderen, een
        grijze ochtend kan bijvoorbeeld een zonnige middag worden.
  \item Waar.
  \item Waar: de Zon is zo fel dat ze het oog voor altijd kan beschadigen,
        zelfs in een ogenblik.
\end{enumerate}
\end{solution}

\begin{solution}{exo:g1:watching-the-weather:8}
Middagen met regen of buien: maandag (regen), vrijdag (buien) en zondag
(regen), dus 3. Koude dagen: maandag, dinsdag en zondag, dus 3. De andere
4 dagen waren warm, dus de week was vooral warm.
\end{solution}

\begin{solution}{exo:g1:watching-the-weather:9}
Ochtend: \omwx{cloud} (grijze lucht); middageten: \omwx{rain} (regen);
avond: \omwx{suncloud} (blauwe lucht met een paar wolkjes).
\end{solution}

\begin{solution}{exo:g1:watching-the-weather:10}
Een vlag wijst naar waar de wind heen gaat. Wijst ze naar links, dan gaat
de wind naar links, dus komt hij van rechts.
\end{solution}

\begin{solution}{exo:g1:watching-the-weather:11}
Geen bladeren: in de winter. Bloesem: in de lente. Bij de winter staat het
sneeuwsymbool: in de winter is het koud, en op veel plaatsen valt er alleen
in dat seizoen sneeuw.
\end{solution}
